\begin{textsample}{POS Dim 1 – mt – Score 74.00 – mt/mt\_em\_30.txt}  \label{ex:f1_pos_072}
6.4 Metodologias para o ensino de \textbf{Ciências} da Natureza e suas \textbf{tecnologias} e as aprendizagens significativas O ensino de \textbf{Ciências} da Natureza, no Ensino Médio, deve transpor a ideia de ensino em um viés disciplinar puro, onde cada componente curricular ocupa um espaço definido no currículo em ação.
Nesse contexto, é importante \textbf{salientar} que, na sociedade \textbf{contemporânea}, temos uma \textbf{juventude} imersa em \textbf{tecnologias} onde as informações são fluidas e complementares, isto é, não mais estáticas e sólidas ( BAUMAN, 2001 ).
Sendo assim, a perspectiva de Bauman corrobora com uma \textbf{abordagem} interdisciplinar para o ensino de Ciências, pois é \textbf{imprescindível} que os estudantes percebam que os conhecimentos das “ Ciências da Natureza e suas Tecnologias ” foram produzidos por meio da ação humana e passam por transformações constantes na sociedade \textbf{contemporânea}.
Nesta \textbf{vertente}, a divisão entre as Ciências da Natureza e Ciências Humanas, por exemplo, não é coerente, pois a Física, a Química e a Biologia são também Ciências Humanas, visto que são conhecimentos estabelecidos pelos humanos ( CHASSOT, 2013 ).
Portanto, uma \textbf{abordagem} interdisciplinar, no espaço da sala de \textbf{aula}, torna -se um ponto indissociável da concepção de formação integral dos estudantes, uma vez que os conhecimentos científicos estão \textbf{imbricados} ao processo de desenvolvimento da \textbf{humanidade}.
Nessa conjuntura, é importante apresentar as Ciências aos estudantes como parte de um contexto histórico, levando em conta que nem sempre os \textbf{erros} na \textbf{ciência} levam a resultados desastrosos, visto que, por exemplo, a borracha galvanizada, o teflon e a penicilina foram resultados de \textbf{erros}.
Além disso, citando caso análogo, Camillo Golgi descobriu a coloração com o ósmio, uma técnica para captar detalhes de neurônios visíveis, depois de espirrar o elemento em um tecido cerebral acidentalmente ( KEAN, 2011 ).
Assim, humanizar os conhecimentos científicos, durante as \textbf{aulas}, é um meio para promover a formação integral dos estudantes, de modo que estes sejam protagonistas do seu projeto de vida, capazes de exercer sua cidadania.
Nessa perspectiva, o ensino de \textbf{Ciências} da Natureza, no contexto da BNCC, não pode ignorar que a sociedade do conhecimento é \textbf{baseada} no desenvolvimento de \textbf{competências} cognitivas.
Isso \textbf{posto}, vale \textbf{frisar} que, para Moran : > A escola padronizada, que ensina e avalia a todos de forma igual e \textbf{exige} resultados previsíveis, ignora que a sociedade do conhecimento é \textbf{baseada} em \textbf{competências} cognitivas, pessoais e sociais, que não se adquirem da forma convencional e que \textbf{exigem} proatividade, colaboração, personalização e visão empreendedora. ( 2015, p.16 ).
Ainda de acordo com Moran ( 2010 ), existe um \textbf{consenso} de que estudantes e professores de que as \textbf{aulas} tradicionais ( expositivas ) pautadas na relação teoria-exercício-teoria, em uma perspectiva de que o conhecimento científico se dá por meio de \textbf{memorização} de fórmulas, fatos e \textbf{teorias}, estão ultrapassadas e precisam de mudanças para que os estudantes possam se sentir \textbf{instigados} ao desenvolvimento de \textbf{competências} e habilidades.
Nesse sentido, Barbosa e Moura ( 2013 ) afirmam que muitos professores utilizam metodologias que podem ser consideradas como tipos de metodologias ativas.
Os professores conhecem, mesmo que de forma inconsciente, meios para propiciar ao estudante um processo de aprendizagem mais ativo.
Neste viés, enfatiza -se que : > As metodologias precisam acompanhar os objetivos pretendidos.
Se queremos que os \textbf{alunos} sejam proativos, precisamos adotar metodologias em que os \textbf{alunos} se envolvam em atividades cada vez mais \textbf{complexas}, em que tenham que tomar decisões e avaliar os resultados, com apoio de materiais \textbf{relevantes}.
Se queremos que sejam criativos, eles precisam experimentar \textbf{inúmeras} novas possibilidades de mostrar sua iniciativa. ( MORAN, 2015, p.17 ).
Dessa forma, cresce a necessidade de se trabalhar com as metodologias em que o professor utilize formas de mediação com o propósito de que os estudantes sejam ( e que \textbf{desperte} o desejo de ser ) protagonistas do processo de aprendizagem.
No contexto da BNCC e do DRC/MT–EM, o processo de ensino e aprendizagem planejado com metodologias ativas é de grande relevância, visto que as \textbf{competências} específicas da área são articuladas para que o estudante, a partir delas, se aproprie das \textbf{competências} gerais preconizadas na BNCC.
Assim, é importante a inclusão de metodologias que possibilitem que os estudantes se tornem mais ativos no processo de ensino e aprendizagem, capazes de compreender e interpretar o mundo ( natural, social e tecnológico ) e de transformá -lo com base nos aportes \textbf{teóricos} e processuais das \textbf{ciências} ( BNCC, 2017 ), ou seja, o desenvolvimento do \textbf{Letramento} Científico nos estudantes torna -se \textbf{imprescindível} para a implementação do DRC/MT–EM.
Nessa perspectiva, considera -se que a literatura tem uma vasta referência de metodologias, \textbf{métodos} ou \textbf{abordagens} didáticas que vão ao encontro do \textbf{exposto}.
É importante \textbf{frisar} que, nesse texto, almeja -se apresentar algumas perspectivas de metodologias, consideradas como ativas, que podem \textbf{auxiliar} os professores na mediação das \textbf{competências} e habilidades das Ciências da Natureza e suas Tecnologias.
Dessa forma, para o desenvolvimento das \textbf{competências} das Ciências da Natureza e suas Tecnologias, é de \textbf{suma} importância que os estudantes consigam desenvolver uma visão integrada da interação existente entre a natureza, o conhecimento científico e o desenvolvimento da \textbf{humanidade} e, portanto, consigam \textbf{mobilizar} os conhecimentos construídos no dia a dia, nas \textbf{aulas} de Física, Química e Biologia, para além da observação dos fenômenos naturais, mas na esfera da intervenção social e aplicação no contexto do mundo do trabalho.
Mediante o apresentado no \textbf{parágrafo} anterior, a nível de exemplo, pode -se afirmar que o estudante, ao ler uma reportagem ou ouvir uma notícia nos veículos de comunicação de que o gás utilizado nos fogões é denominado de gás GLP ( Gás Liquefeito do Petróleo ), precisa ter consciência de que o gás GLP é composto de uma mistura de hidrocarbonetos de baixo peso molecular.
Além disso, precisa, sempre que necessário, ser capaz de interpretar estas informações em diferentes contextos, por exemplo, na perspectiva dos processos que ocorreram na natureza ( os processos realizados pelas bactérias, transporte de energia e calor nos sistemas ambientais ) para a formação do petróleo, e também as \textbf{tecnologias} envolvidas em sua extração para produção de combustíveis.
Nessa \textbf{vertente}, para formar sujeitos capazes de utilizar os conhecimentos das Ciências da Natureza e suas Tecnologias de modo \textbf{crítico} e integrado, o Ensino por Investigação apresenta -se como uma possibilidade para o professor no processo de mediação das habilidades da \textbf{competência} específica das Ciências da Natureza e suas Tecnologias.
Assim, por meio do Ensino por Investigação, o estudante será \textbf{instigado} a \textbf{argumentar}, apresentar hipóteses, propor a resolução do \textbf{problema} em questão, analisar os dados e comunicar os resultados.
O professor é peça-chave para favorecer o ensino de \textbf{maneira} que o estudante se torne mais ativo.
Porém, o trabalho do dia a dia com metodologias e estratégias diferenciadas de ensino não é \textbf{tarefa} fácil, mas são essenciais.
Por isso, é preciso pensar em um ensino que proporcione ao estudante um maior protagonismo, seja com materiais de baixo custo, ou feito por meio de uma demonstração investigativa, que ocorre, por exemplo, quando o professor não tem acesso a um lugar específico, como laboratório, e a atividade experimental é realizada em sala de \textbf{aula}.
Ao professor cabe o questionamento, a inclusão de \textbf{problematizações} ou perguntas que façam com que os estudantes possam pensar em uma \textbf{maneira} de como resolver o \textbf{problema} proposto.
A \textbf{maneira} de conduzir a \textbf{aula} é diferenciada, porque o professor faz as perguntas, mas não dá as respostas de \textbf{imediato}.
Isso faz com que os estudantes, além de saírem de suas áreas de conforto, também sejam desafiados a resolverem o \textbf{problema}.
Portanto, o Ensino por Investigação trabalha com a \textbf{vertente} de que o estudante se envolva e se comprometa, tornando -se mais ativo em todo o processo.
Além do Ensino por Investigação, o Estudo de Caso \textbf{configura} -se como um meio para \textbf{auxiliar} os estudantes a fundamentar e defender decisões éticas e responsáveis de modo colaborativo e com o viés interdisciplinar, corroborando para o processo de construção das \textbf{competências} específicas 1, 2 e 3.
Sá e Queiroz ( 2010 ) enfatizam que o \textbf{método} de Estudo de Caso é uma variante do Aprendizado \textbf{Baseado} em \textbf{Problemas} ( ABP ).
Neste sentido, destaca -se que o \textbf{método} de Estudo de Caso possibilita ao professor trabalhar com \textbf{problemas} reais, com o propósito de estimular o desenvolvimento do pensamento \textbf{crítico}, a habilidade de resolução de \textbf{problemas} e a aprendizagem de conceitos ( SÁ E QUEIROZ, 2010 ).
Sendo assim, o \textbf{método} impulsiona os estudantes a buscarem respostas e tomarem decisão sobre a melhor forma de resolver \textbf{problemas} \textbf{pertinentes} para a sociedade.
Com a perspectiva de que o estudante se comporte diferentemente do que está acostumado, \textbf{tanto} o Ensino por Investigação quanto o Estudo de Caso trazem o \textbf{problema} a ser resolvido como ponto-chave para que o estudante inicie um processo de investigação e passe a se familiarizar com normas e valores da comunidade científica.
Assim se evitará que se construa visão deformada sobre \textbf{ciência}, assim como o trabalho do \textbf{cientista}.
Portanto, estas metodologias apresentam -se como possibilidades para serem utilizadas no processo de mediação dos professores para \textbf{desencadear} processos que viabilizem a apropriação das \textbf{competências} específicas das Ciências da Natureza e também as \textbf{competências} gerais da BNCC.
Outro processo \textbf{metodológico} que pode ser amplamente utilizado pelos professores para \textbf{instigar} o processo de investigação científica nos estudantes é a Aprendizagem \textbf{Baseada} em \textbf{Problemas} ( ABP ), uma estratégia de mediação de conhecimentos que se organiza ao redor da investigação de \textbf{problemas} do mundo real ( LOPES et al., 2019 ).
Estudantes e professores se envolvem em analisar, compreender e propor soluções para situações cuidadosamente desenhadas de modo a garantir ao estudante a apropriação das \textbf{competências} específicas 1, 2 e 3, previstas no DRC/MT–EM.
Desta forma, a ABP possibilita que o estudante aplique seus conhecimentos na elaboração de \textbf{argumentos} e \textbf{posicionamentos} frente aos diferentes cenários de \textbf{problematização}, de modo colaborativo e com ética e responsabilidade \textbf{socioambiental}.
Nesse processo de elaboração de \textbf{argumentos} no âmbito da ABP, enfatiza -se a possibilidade de, por meio da mediação do professor, o estudante desenvolver a visão da \textbf{ciência} como um processo contínuo e progressivo \textbf{inerente} às necessidades humanas.
Nessa perspectiva, a ABP é um importante ponto de partida para o planejamento dos professores da área, principalmente para a mediação da \textbf{competência} específica 3, visto que proporciona aos professores meios para motivar os estudantes na \textbf{mobilização} dos conhecimentos científicos na resolução de situações-problema, no contexto de uma visão mais interdisciplinar e transdisciplinar.
Além disso, em síntese, pode -se resumir que a ABP é uma metodologia que se insere no contexto da BNCC como uma importante estratégia de mediação para desenvolver o protagonismo juvenil e \textbf{instigar} o uso de \textbf{tecnologias} \textbf{digitais} de informação e comunicação como meios para o processo de investigação científica.
No Quadro 5, apresenta -se uma síntese das etapas que os professores precisam pensar durante a mediação por meio da ABP. \textbf{Quadro} 5 - Os “ Sete Saltos ” ( SCHMIDT, 1983 ; FINCO-MAIDAME ; MESQUITA, 2017. ) 1.
Esclarecer frases e conceitos confusos na formulação do \textbf{problema}. 2.
Definir o \textbf{problema} : descrever exatamente que fenômenos devem ser entendidos e explicados. 3.
Chuva de ideias ( Brainstorming ) : usar conhecimentos prévios e senso comum.
Tentar formular o máximo possível de explicações. 4. \textbf{Detalhar} as explicações propostas : tentar construir uma “ \textbf{teoria} ” pessoal, coerente e \textbf{detalhada} dos processos subjacentes aos fenômenos. 5.
Propor temas para a aprendizagem autodirigida. 6.
Preencher as lacunas do conhecimento por meio do estudo individual. 7.
Compartilhar as conclusões com o grupo e integrar os conhecimentos adquiridos em uma explicação adequada dos fenômenos.
Comprovar o que sabe.
Avaliar o processo de construção de conhecimentos.
Dessa forma, é de \textbf{primordial} necessidade o olhar atento do professor, que atuará como mediador ou orientador do processo, considerando o senso comum e os conhecimentos prévios para o desenvolvimento dos novos conhecimentos.
Nesse sentido, a construção do conhecimento na escola \textbf{depende} intrinsecamente da relação de mediação pedagógica estabelecida pelo professor, pois, se o objetivo é formar um cidadão com o desenvolvimento das dez \textbf{competências} gerais, apresentadas na BNCC, é de fundamental importância que a linguagem científica, a criatividade e o protagonismo sejam elementos estruturantes do desenvolvimento do estudante no espaço escolar.
Neste contexto, é \textbf{relevante} reforçar que a ABP é uma metodologia que apresenta variantes, sendo que, para além do Estudo de Caso, conforme já enfatizado em \textbf{parágrafos} anteriores, a aprendizagem \textbf{baseada} em projetos também é uma metodologia variante da ABP.
De forma geral, a aprendizagem \textbf{baseada} em projetos, segundo Bender ( 2014 ), pode ser definida pela utilização de projetos \textbf{autênticos} e realistas, \textbf{baseados} em uma questão, \textbf{tarefa} ou \textbf{problema} que \textbf{instigue} um processo de aprendizagem significativa.
Em síntese, os principais objetivos da aprendizagem \textbf{baseada} em projetos são estimular a motivação dos estudantes, promover uma aprendizagem centrada, incentivar o trabalho em grupo, desenvolver o espírito de iniciativa e criatividade, \textbf{aprimorar} \textbf{capacidades} de comunicação, potencializar o pensamento \textbf{crítico} e proporcionar estratégias para que as habilidades das Ciências da Natureza e suas Tecnologias sejam trabalhadas de forma interdisciplinar ( ALVES ; MOREIRA ; SOUZA, 2007 ). \textbf{Posto} isto, no contexto de proporcionar o desenvolvimento das dez \textbf{competências} gerais apresentadas na BNCC, é de \textbf{suma} relevância articular a aprendizagem \textbf{baseada} em projetos e a \textbf{abordagem} STEAM ( Ciência, Tecnologia, Engenharias, Artes e Matemática, na sigla em \textbf{inglês} – \textbf{Figura} 4 ).
De forma geral, observa -se, por meio da \textbf{Figura} 4, que a \textbf{abordagem} STEAM pode proporcionar aos professores de \textbf{ciências} um planejamento pedagógico mais integrado com as \textbf{demandas} da sociedade \textbf{contemporânea}.
Segundo Bacich e Holanda ( 2020 ), se essa \textbf{abordagem} for associada à metodologia de aprendizagem \textbf{baseada} em projetos, os estudantes poderão ampliar o senso de relevância dos conhecimentos científicos desenvolvidos na educação básica.
Em linhas gerais, alguns elementos são comuns para a aprendizagem \textbf{baseada} em projetos e, consequentemente, também os articuladores dos projetos de STEAM, que são : > [... ] um contexto \textbf{autêntico} capaz de engajar os estudantes, de preferência relacionado a um \textbf{problema} real ; uma \textbf{sequência} de etapas organizadas para a exploração do conhecimento científico e para a produção dos estudantes ; um produto \textbf{final}, geralmente um artefato que permita a aplicação das ideias das engenharias ; e, por fim, a comunicação do projeto, para compartilhar com a comunidade e \textbf{sistematizar} suas aprendizagens. ( BACICH e HOLANDA, 2020, p.19 ) Pensando sob essa forma de organizar o planejamento dos professores de \textbf{ciências}, a \textbf{abordagem} em STEAM, com a metodologia aprendizagem \textbf{baseada} em projetos, precisa ser compreendida como um meio de estimular a criatividade, a investigação científica e a cooperação entre os pares.
Assim, por exemplo, ao propor que os estudantes avaliem o impacto ambiental do descarte inadequado de pneus nas adjacências da escola, e posteriormente solicitar que desenvolvam uma estratégia de intervenção social e ambiental, o fato de os professores da área das Ciências da Natureza debaterem o processo de fabricação dos pneus e também as possíveis doenças tropicais associadas ao descarte inadequado destes, com os estudantes, não se \textbf{configura} como uma \textbf{abordagem} em STEAM ; e os professores de Arte direcionarem os estudantes para a produção de um jardim na escola, a partir da reciclagem dos pneus, pintando -os, como forma de intervenção \textbf{socioambiental}, também não. \textbf{Convém} destacar que o meio de desenvolvimento de um projeto, conforme supracitado, não é inadequado.
Porém, se o objetivo for trabalhar com os estudantes por meio de uma \textbf{abordagem} em STEAM, as ações elencadas necessitam estabelecer uma relação \textbf{dialógica} com os princípios articuladores da \textbf{abordagem} STEAM, visto que, para promovê -la, é \textbf{imprescindível} que os estudantes sejam \textbf{instigados} a analisarem \textbf{problemas} \textbf{complexos} que não possuem uma única solução, mas que possuam contexto \textbf{autêntico}, com viés tecnológico para além do mundo \textbf{digital}, ou seja, como forma de conectar conhecimentos científicos para o desenvolvimento de produtos ( BACICH e HOLANDA, 2020 ).
Além disso, vale \textbf{frisar} que, para formar estudantes aptos à tomada de decisões no campo profissional e pessoal, é importante uma postura de mediação por parte dos professores em um projeto STEAM, ou seja, não direcionar como será o processo de intervenção dos estudantes, mas \textbf{despertar} esses processos criativos e atitudes.
Portanto, a respeito do exemplo mencionado anteriormente, há que se motivar os estudantes a um processo \textbf{reflexivo} em que a \textbf{compreensão} dos conhecimentos científicos sobre as propriedades dos materiais, a \textbf{análise} dos processos de reciclagem e a necessidade de garantir a saúde coletiva sejam uma constante.
Estes são pontos de atenção no planejamento pedagógico dos professores, uma vez que é de \textbf{suma} importância que sejam evitadas, durante a realização do projeto, a \textbf{simples} aplicação de conceitos científicos para a resolução de um \textbf{problema} \textbf{complexo} ou aplicação de roteiros para produção de artefatos.
Nesta perspectiva, enfatiza -se também a necessidade de garantir o alinhamento das \textbf{competências} e habilidades apresentadas neste documento como uma intencionalidade pedagógica do projeto STEAM.
Por último, \textbf{convém} \textbf{salientar} que as proposições dos \textbf{métodos} de ensino, supramencionados, são possibilidades que devem ser expandidas, no cotidiano das escolas, a partir de uma prática pedagógica \textbf{reflexiva} e \textbf{crítica} dos professores das Ciências da Natureza e suas Tecnologias.

% matched lemmas: abordagem, aluno, análise, aprimorar, argumentar, argumento, aula, autêntico, auxiliar, basear, capacidade, cientista, ciência, competência, complexo, compreensão, configurar, consenso, contemporâneo, convir, crítico, demanda, depender, desencadear, despertar, detalhar, dialógico, digital, erro, exigir, exposto, figura, final, frisar, humanidade, imbricar, imediato, imprescindível, inerente, inglês, instigar, inúmero, juventude, letramento, maneira, memorização, metodológico, mobilizar, mobilização, método, parágrafo, pertinente, posicionamento, postar, primordial, problema, problematização, quadro, reflexivo, relevante, salientar, sequência, simples, sistematizar, socioambiental, sumo, tanto, tarefa, tecnologia, teoria, teórico, vertente
\end{textsample}
